% MAS1004 Assignment 1: long report.
%
% You fill this in with your agent, in your repository. AGENTS.md tells the
% agent how. Red text in the PDF is what still has to be replaced.
%
% - Pictures go in figures/. Copy them from results/, which is not in git.
% - Outputs go in outputs/. run.py saves what it prints in
%   results/<tag>_output.txt, and PLAN.md saves the other outputs in results/
%   with `| tee results/...`. Copy them here. \outputfile{outputs/check.txt}
%   prints a file exactly as it is. Every number in this report has to come
%   from such an output.
%
% To make the PDF: compress this folder into a zip, and on Overleaf choose
% New Project, Upload Project. Set Menu, Compiler to XeLaTeX, and Recompile.

\documentclass[11pt]{article}
\usepackage{mas1004}

\reporttitle{Assignment 1: Long report}
\studentname{Abukar Hassan}
\studentid{G20260557}

\begin{document}
\makeheader

% ---------------------------------------------------------------------------
\section*{Problem 1. Categories and images}

\subsection*{Why these categories}
For my three categories I picked toys, specifically a Trafalgar Law Lego piece,
a sleepy Kirby, and a Toothless keychain. The reason why I picked these toys is
because they all hold meaning to me. For instance, I picked the Trafalgar Law
Lego piece because he's my favorite character from the Japanese anime
``One Piece''. The same could be said for the rest of the toys. They are my
favorite characters from their respective fields of media. Another reason why I
picked them is because I recently acquired them during my time here in Korea. I
took 10 pictures of each toy from different angles and with different
backgrounds to see how the model will differentiate them.

\subsection*{What I expected the model to use}
One way I believe the model will be able to tell them apart is by using the
color of the characters because they all have different color ways. As Kirby is
mostly pink, Toothless is black, Trafalgar Law has a hint of blue which none of
the others have. Another way it could tell them apart would be by the faces
since each of the toys have vastly different facial features. Lastly, the model
could tell them apart by the shape of each of the characters since they all
look vastly different.

\subsection*{Images per category}
Image counts: Kirby: 150 images found. Toothless: 150 images found. Trafalgar
Law: 150 images found.

\outputfile{outputs/collect_step2.txt}

% ---------------------------------------------------------------------------
\section*{Problem 2. Preparing an image}

\sidebyside{figures/prepare_original.jpg}{the original photo, IMG\_2645}
           {figures/prepare_result.png}{what prepare\_image made of it}

trafalgar\_law\_figure\_IMG\_2645 shows a tall image of my Law LEGO figure.
Prepare\_image resized the short side of the image to 256, which meant that the
long side was also scaled with it. This caused Law's hat, part of his face, and
legs to be cut off due to the photo being too tall. This could cause some
problems because his hat is iconic and his face is different from the other
toys. These could be what the model uses to differentiate the toys. Now the
model can't use his hat, half his face or legs to figure out which toy it is
based on.

% ---------------------------------------------------------------------------
\section*{Problem 3. First training runs}

\begin{center}\footnotesize
\begin{tabular}{llllrrrr}
\toprule
tag & start & trained & epochs & lr & trainable parameters & train acc. & test acc. \\
\midrule
run & ImageNet & all layers & 10 & 0.0001 & 11,178,051 & 100.0\% & 100.0\% \\
scratch & random & all layers & 10 & 0.0001 & 11,178,051 & 100.0\% & 93.3\% \\
frozen & ImageNet & last layer & 10 & 0.001 & 1,539 & 100.0\% & 97.8\% \\
bigLR & ImageNet & all layers & 10 & 0.01 & 11,178,051 & 68.1\% & 70.0\% \\
\bottomrule
\end{tabular}
\end{center}

\picturefile{figures/run_curves.png}

The loss went down sharply during training but then flattened closer to the
0.00 loss range when it reached 3-4 epochs. Additionally, while the training
line started high the actual test line started much lower and closer to where
the training line started to flatten out.

Interestingly enough for the first training run there wasn't a gap between the
training (100\%) and the test (100\%). This means that there wasn't a big
difference between the two. For instance many of the photos that the web search
generated had a couple duplicates and pictures from the same angle which could
have contributed to the high test result.

\begin{center}\picturefile[0.5\linewidth]{figures/run_confusion.png}\end{center}

For the main run's confusion matrix the model did not mix any of the
categories it got 100\% on both the training and test. This did not surprise me
since the categories look pretty different from each other. So I looked at
random start and found some mistakes, for example during the random test the
model had 3 mistakes where it confused Toothless for Law. This surprised me
because they do not look alike at all. I think that the model confused the two
because most of the pictures for both images have a white background and the
biggest thing that I believe contributed to the confusion is Law's clothing
color which the model could have thought was Toothless because both are black.

\begin{center}\picturefile[0.5\linewidth]{figures/scratch_confusion.png}\end{center}

The difference between the ImageNet (100\%) and the Random (93.3\%) test scores
is 6.7\%. I think the reason why the ImageNet test score was higher is because
it learned from many other images that helped it recognize things like eyes,
wings, capes, shapes, and more. This gave it a better chance of scoring higher
than the random test because the random model does not have prior knowledge
which could have helped in differentiating the toys. Additionally the random
model only had 360 images to train with and could have learned to recognize the
differences in the toys. Another reason could be that that model memorised the
differences.

The end result of training only the last layer was that it scored 2.2\% lower
(97.8\% vs 100\%), by only training 1,539 of the total 11,178,051 parameters.

\subsection*{Output of the runs}
\outputfile{outputs/run_output.txt}
\outputfile{outputs/scratch_output.txt}
\outputfile{outputs/frozen_output.txt}
\outputfile{outputs/bigLR_output.txt}

% ---------------------------------------------------------------------------
\section*{Problem 4. Cleaning}

\subsection*{My rule}
Rule: toys that are in the box only, have other characters alongside my
character (including human faces), wrong characters all together, and near
duplicates are to be removed. Toys that are outside or near the box, only my
character (or multiple images of my character), and human hands alone are okay
to keep.

\subsection*{Three kinds of junk}
Packages / boxes (kirby\_figure/0001.jpg):
\picturefile[0.3\linewidth]{figures/p4_kirby_0001_box.png}

Many characters in the image aside from the selected character
(kirby\_figure/0124.jpg):
\picturefile[0.3\linewidth]{figures/p4_kirby_0124_many_characters.png}

Wrong character altogether (kirby\_figure/0137.jpg):
\picturefile[0.3\linewidth]{figures/p4_kirby_0137_wrong.png}

\subsection*{What I removed}
\outputfile{outputs/clean_count.txt}

Near copies found: I found 48 near copies. I first removed the ones that were
counted as ``Junk'' and then kept exactly one from each of the duplicates.

Out of the first 20 suspects for each category, I removed 8 from Kirby, 15
from Toothless, and 5 from Law. I missed one of the images for junk when it
came to Law because I didn't see the leopard under him.

\outputfile{outputs/suspects_removed.txt}

Yes I found junk that the tool didn't show that went against my packaging rule.
For instance, kirby\_figure/0001.jpg which I showed as a junk example for
packages/boxes.

\subsection*{Before and after cleaning}
\outputfile{outputs/compare.txt}

When comparing the numbers for before and after the cleaning the number did
not change it stayed the same with 76.7\% for the run when compared to the
clean score of 76.7\%. The reason why I think the number did not move is
because the junk I deleted was not something that the model used when
differentiating my toys. A reason why I removed the boxes is because my toys do
not have boxes in them.

% ---------------------------------------------------------------------------
\section*{Problem 5. Images from a new source}

My images came from me, I took these pictures in my dorm room of toys I own.
The reason why I'm sure they are not among the downloaded images is because I
personally took these pictures and I reviewed all the pictures that were
downloaded.

\outputfile{outputs/overlap.txt}

The accuracy on my downloaded test set and on my new images:

\begin{center}\small
\begin{tabular}{llr}
\toprule
 & downloaded & my new images \\
\midrule
clean & 100\% & 76.7\% \\
\bottomrule
\end{tabular}
\end{center}

5 mistakes it made: Before stating the mistakes I want to say that only the
first 2 images are confident mistakes of 70\% or higher while the rest are
around the 50\% mark. As for the rotated images I want to point out again that
the model got 3/10 when the images were stored sideways but 7/10 when retested
upright.

1. IMG\_2648.jpg (Guessed Toothless but was Law). The thing that pushed it
away from Trafalgar Law is that this image wasn't showing his face.

\picturefile[0.3\linewidth]{figures/p5_2648.png}

2. IMG\_2650.jpg (Guessed Toothless but was Law).

\sidebyside{figures/p5_2650_sideways.png}{what the model saw}
           {figures/p5_2650_upright.png}{the image as it was taken}

The reason why the model got them wrong was because the model saw the image
sideways instead of its intended upright position.

3. IMG\_2652.jpg (Guessed Toothless but was Law).

\sidebyside{figures/p5_2652_sideways.png}{what the model saw}
           {figures/p5_2652_upright.png}{the image as it was taken}

The reason why the model got them wrong is exactly the same reason as number 2.

4. IMG\_2645.jpg (Guessed Toothless but was Law). This image only shows part
of his face, it could also be the color of his clothing.

\picturefile[0.3\linewidth]{figures/p5_2645.png}

5. IMG\_2701.jpg (Guessed Toothless but was Law).

\sidebyside{figures/p5_2701_sideways.png}{what the model saw}
           {figures/p5_2701_upright.png}{what the intended image should have
           looked like}

The reason why the model got them wrong is exactly the same reason as number 2
and 3.

The main difference about my new images is that some of them were stored
sideways on accident. Other things that I did to make my images different from
the downloaded ones was take pictures from different angles and distances.
Lastly, I took pictures with different backgrounds such as tables, near
windows, on the bed, and with a dark background.

% ---------------------------------------------------------------------------
\section*{Problem 6. The demo}

Address: \url{https://abukarhassan2222.github.io/mas1004-assignment1/}

I gave it a picture of a glass of water with a white background and the
answers it gave were all over the place with the highest percentage being
50.2\% for Kirby.

% ---------------------------------------------------------------------------
\section*{Problem 7. What the model looks at}

% Copy this subsection once for every guess.
\subsection*{Guess 1}

My guess: It tells a Trafalgar Law image from a Toothless or Kirby by his
face/hat. Reasoning: from what I saw in the demo, Law's face is what the model
uses to figure out it's him.

I would like to cover the images where Law's face shows to check if my guess
is correct. This change will test my guess is because the model will need to
use something else aside from his face to guess what the image is showing.

My request to the agent, word for word:
\begin{quote}
make a copy of each image in data/new\_images/trafalgar\_law\_figure with the
top third covered by a grey rectangle , and save the copies in
data/changed/trafalgar\_law\_figure\_no\_head
\end{quote}

I noticed that the agent made a mistake with the copies that had the change in
them, where it rotated them sideways. This could be because some of the photos
were taken like that. So I gave the agent a new prompt.

Prompt:
\begin{quote}
the copies are wrong beacuse they were stored sideways so can you trun them
upright like the orginals. also cover a little bit more than a third of the
image and leave out photots where it doesnt cover the face
\end{quote}

After this prompt it did fix all the images that were sideways. But there were
still some images that still had his face on it. So I gave cline another
prompt to help with that.

Prompt 3:
\begin{quote}
there are still some images that still has his face in them. IMG\_2650.jpg,
IMG\_2651.jpg, IMG\_2652.jpg, IMG\_2653.jpg, IMG\_2695.jpg , IMG\_2701.jpg.
can you cover another 10\% of each of these images with the grey rectangle but
leave the rest untouched
\end{quote}

Now most of the images had his face covered except 2 which were images
IMG\_2651.jpg and IMG\_2652.jpg. So I gave it another prompt.

Prompt 4:
\begin{quote}
there are now 2 images left that still need a little more covering which are
IMG\_2651.jpg and IMG\_2652.jpg. cover another 5\% of these images while
leaving the rest as they are
\end{quote}

This worked and finally I have all the faces covered for law.

While working on the demo I noticed that my numbers from question 5 did not
add up to what I was seeing on the demo at all. In step 15, I found that the
model only got 3/10 for law while my demo showed that it actually got 8/10
right. The only thing that changed was that I rotated the copied images and
retested them. This could mean that during step 15 the images that were
guessed wrong were actually rotated sideways instead of being right up which
could have contributed to such a low score.

New prompt:
\begin{quote}
something isnt adding up when i ran the numbers with the demo compared to what
I got for step 15. it turns out that the model actually gets 8/10 right when
it comes to trafalgar law but earlier it showed 3/10. The only thing that was
changed was the rotation of the images. so im wondering if prepare\_image
turns the photos upright or leaves it as? and which photos were stored
sideways. ONLY CHECK DONT CHANGE ANYTHING
\end{quote}

After that I prompted it again with:
\begin{quote}
ill keep the data and explain the rotation discrepancy in the report. can you
save /tmp/check\_orientation\_test.py as a file in the and adding it into
mas1004-assignment1. also save this test into this file too so it can be re
run later so my professor can test it and get the same results
\end{quote}

\outputfile{outputs/check_orientation.txt}

My conclusion. My initial guess showed that it's very likely that the model
uses Law's face to tell him apart from Kirby and Toothless. I am pretty sure
of it since the model went from 8/10 correct all the way down to 3/10 when I
tested it on the demo.

Law face vs no face:

\begin{center}\footnotesize
\begin{tabular}{l p{0.30\linewidth} p{0.42\linewidth}}
\toprule
Image & Model's confidence before change & Model's confidence after change \\
\midrule
IMG\_2644.jpg & 78.5\% law (model actually sees a little more than half his face) & 60.7\% law \\
IMG\_2645.jpg & 48.5\% law (model actually sees half his face from an angle) & 56.4\% toothless and 32.1\% Law \\
IMG\_2647.jpg & 37.7\% kirby, 31.3\% Law, and 31\% toothless (over head image with his round hat in the middle) & 59.9\% toothless, 25.9\% Law. The jump in toothless score could be due to the face/hat being covered and only relying on the color of the image which is similar to toothless because he is also black and blue in some images) \\
IMG\_2648.jpg & 67.8\% toothless and 27.4\% law (this image is from behind so his face is not shown but his hair can be seen a little. Most of the image is his clothing which is again black and blue) & Interesting enough, toothless score went down to 52.7\% and law went down to 26.4 \\
IMG\_2650.jpg & 86.3\% law (The image shows a higher angle but you can still see a small part of laws face) & Now the model thinks its kirby with a 51.2\% score while law has 13.5\% \\
IMG\_2651.jpg & 90.3\% law (image is similar to the last one but it shows a little more of his face in this one) & 41.8\% for kirby and 37.3\% for law \\
IMG\_2652.jpg & 99.1\% that its is law (the model can see laws full face) & With his face cut off the model shows a 56.3\% confidence its law \\
IMG\_2653.jpg & 89.3\% that its law (model can see laws side profile including his hair and a little bit of his eye) & Model with his face cut off now shows a 48.1\% confidence that its law \\
IMG\_2695.jpg & 95\% that its law (image sees his whole face) & With law's face cut off the model now shows a 40.5\% confidence that its toothless which was previously 0.9\% and law went down to 25.6\% \\
IMG\_2701.jpg & 96.4\% that its law (model can see his whole face) & With his face cut off the model thinks its an image of kirby with 64.3\% confidence but its actually law \\
\bottomrule
\end{tabular}
\end{center}

Example 1: model is very confident that the image is 99.1\% its law but when I block his face the model loses 42.8\% confidence because
of that change. Showing just how much that change affects the models choices.

\sidebyside{figures/p7_2652_before.png}{original: 99.1\% Law}
           {figures/p7_2652_after.png}{changed: 56.3\% Law}

Example 2: the model has a 95\% confidence that the model shows Law. But, when
I cover his face the model thinks its Toothless with a 40.5\% confidence and
only a 25.6\% confidence for Law.

\sidebyside{figures/p7_2695_before.png}{original: 95\% Law}
           {figures/p7_2695_after.png}{changed: Toothless 40.5\%, Law 25.6\%}

Example 3: the model shows a 96.4\% confidence in it being law. But, after
covering his face the model shows a 64.3\% confidence that its Kirby instead.

\sidebyside{figures/p7_2701_before.png}{original: 96.4\% Law}
           {figures/p7_2701_after.png}{changed: Kirby 64.3\%}

\sidebyside{figures/p7_2651_before.png}{original: 90.3\% Law}
           {figures/p7_2651_after.png}{changed: Kirby 41.8\%, Law 37.3\%}

\subsection*{Guess 2}

I'd like to come back to my guess from problem one which was the color of each
toy. Since Kirby and Toothless have very different colors I want to see if the
model uses the colors of the images to tell them about.

My request to the agent, word for word:
\begin{quote}
make a copy of each image in data/new\_images/toothless\_dragon\_toy and
data/new\_images/kirby\_figure and turn the images grayscale , and save the
copies in data/changed/kirby\_figure\_gray\_scale and
data/changed/toothless\_dragon\_toy\_gray\_scale
\end{quote}

The model proved my guess wrong as it correctly guessed all the grayscale
images of Kirby and Toothless.

Kirby gray scaling test:

\begin{center}\footnotesize
\begin{tabular}{lrr}
\toprule
Image & Before & After \\
\midrule
IMG\_2654.jpg & 99.8\% kirby & 92.3\% kirby \\
IMG\_2655.jpg & 98.5\% kirby & 89.8\% kirby \\
IMG\_2656.jpg & 98.4\% kirby & 90.2\% kirby \\
IMG\_2657.jpg & 99\% kirby & 97\% kirby \\
IMG\_2658.jpg & 99.3\% kirby & 96.9\% kirby \\
IMG\_2659.jpg & 94\% kirby & 78\% kirby \\
IMG\_2660.jpg & 98.8\% kirby & 91.9\% kirby \\
IMG\_2661.jpg & 98.5\% kirby & 86.3\% kirby \\
IMG\_2694.jpg & 90.7\% kirby & 71.2\% kirby \\
IMG\_2700.jpg & 91.1\% kirby & 69.6\% kirby \\
\bottomrule
\end{tabular}
\end{center}

toothless gray scaling test:

\begin{center}\footnotesize
\begin{tabular}{lrr}
\toprule
Image & Before & After \\
\midrule
IMG\_2662.jpg & 100\% toothless & 100\% toothless \\
IMG\_2664.jpg & 99.9\% toothless & 99.9\% toothless \\
IMG\_2665.jpg & 99.7\% toothless & 99.5\% toothless \\
IMG\_2666.jpg & 100\% toothless & 100\% toothless \\
IMG\_2667.jpg & 98\% toothless & 98.8\% toothless \\
IMG\_2668.jpg & 100\% toothless & 100\% toothless \\
IMG\_2669.jpg & 99.9\% toothless & 99.8\% toothless \\
IMG\_2670.jpg & 93.8\% toothless & 97.6\% toothless \\
IMG\_2696.jpg & 99.8\% toothless & 99.9\% toothless \\
IMG\_2699.jpg & 96.3\% toothless & 96.9\% toothless \\
\bottomrule
\end{tabular}
\end{center}

\subsection*{Guess 3}

I now want to try grayscaling Law because some of the images it confused with
Kirby/Toothless could be due to his clothing (like his hat or pants).

My request to the agent, word for word:
\begin{quote}
can you do the same grayscale to law as well? by making a copy of
data/new\_images/trafalgar\_law\_figure and trun those images grayscale. then
save the copies into data/changed/trafalgar\_law\_figure\_gray\_scale
\end{quote}

Law grayscale demo:

\begin{center}\footnotesize
\begin{tabular}{l p{0.30\linewidth} p{0.40\linewidth}}
\toprule
Image & Model's confidence before change & Model's confidence after change \\
\midrule
IMG\_2644.jpg & 78.5\% law (model actually sees a little more than half his face) & 51.4\% toothless and law at 40.1\% \\
IMG\_2645.jpg & 48.5\% law (model actually sees half his face from an angle) & 70.4\% toothless and 20.6\% Law \\
IMG\_2647.jpg & 37.7\% kirby, 31.3\% Law, and 31\% toothless (over head image with his round hat in the middle) & 60\% toothless, 25.1\% Law. The jump in toothless score could be because of the grayscaling \\
IMG\_2648.jpg & 67.8\% toothless and 27.4\% law (this image is from behind so his face is not shown but his hair can be seen a little. Most of the image is his clothing which is again black and blue) & Interesting enough, toothless score went up to 93.5\% \\
IMG\_2650.jpg & 86.3\% law (The image shows a higher angle but you can still see a small part of laws face) & Now the model thinks its law with a 54.4\% score \\
IMG\_2651.jpg & 90.3\% law (image is similar to the last one but it shows a little more of his face in this one) & 83.4\% for law \\
IMG\_2652.jpg & 99.1\% that its is law (the model can see laws full face) & Down to 96.3\% Law \\
IMG\_2653.jpg & 89.3\% that its law (model can see laws side profile including his hair and a little bit of his eye) & 86.2\% law \\
IMG\_2695.jpg & 95\% it is law (image sees his whole face) & 92.7\% law \\
IMG\_2701.jpg & 96.4\% that it is law (model can see his whole face) & 93.0\% law \\
\bottomrule
\end{tabular}
\end{center}

\subsection*{What I now think the model uses}

I believe the model uses face and shape features to determine which toy is
which while color helps when the model's face is partially visible. I am
pretty sure because when I tested my hypothesis about the model using facial
features I used my law images which was the toy the model had the most trouble
with compared to the other toys (Toothless and Kirby) where the model got
100\% accuracy on. Which leads me to believe that there is something in the
images of Law that was confusing the model. The differences I found between
Law and the other figures was that Law had a hat on which partially covered
his face in certain angles and that his color was slightly less distinct from
Kirby who is almost fully pink/red and toothless who is fully black. When I
tested my hypothesis by putting a grey box over law's face the model guessed
3/10 correctly as law which is different from when the model got 8/10 (demo
test) with his face showing. I then tested my other hypothesis about color in
which the model got 100\% again on toothless and Kirby, but not law. I looked
through the images and found that the images that partially showed his face
but not the whole thing were confused as toothless and not Kirby which could
be due to grayscaling turning Law black just like toothless. But, there were
some results from the demo where I cut law's face off where the model thought
it was Kirby. This could be due to the shape of the toy instead of either the
color or the faces.

Finally, I would like to talk about some findings I found during problem 7.
When testing out if the demo worked correctly I noticed that the numbers I was
getting and the numbers from my first test were different even when I used the
same images. This led me to finding out that my original score of 3/10 was
largely due to the sideways images. I then retested the images and got 7/10 on
the terminal and 8/10 on the demo. This is because python and the web browser
resize images differently which changed the outcome of one image. This caused the inconsistency in my results between the terminal and the demo. The face covered demo shows 3/10 answers correct while the changed copies under check orientation show a result of 2/10.  I decided to
keep those numbers and talk about my findings and mishaps. Nevertheless, this
did not affect the models answers during my demos since all the images were
upright throughout the demos.

% ---------------------------------------------------------------------------
\section*{Appendix. Output of check.py}
\outputfile{outputs/check.txt}

\end{document}
